\begin{scanpage}{0113}{96}
(16.19) by the expression
\[\left.\begin{aligned}E_x=-H_y&=-\frac a2E_{0x}\left(\operatorname{tg}\omega+\frac{\frac2n\sin\frac\pi n}{\cos\frac\pi n-\cos\frac{2\omega}n}\right)\frac{e^{ikR}}R e^{2ikl},\\E_y=H_x&=0,
\end{aligned}\right\}\tag{17.03}\]
where
\[n=1+\frac{\omega+\Omega}\pi.\tag{17.04}\]
An asymptotic calculation of the rigorous diffraction series for a semi-infinite cone [38--40] shows that in the direction $\vartheta=\pi$ one may neglect the effect of the nonuniform part of the current caused by the conical point. Therefore, summing (17.02) and (17.03), we obtain the following expression for the fringing field:
\[\left.\begin{aligned}E_x=-H_y&=-E_{0x}\left[\frac i2\operatorname{tg}^2\omega(1-e^{2ikl})\right.\\&\left.\qquad+ka\frac{\frac2n\sin\frac\pi n}{\cos\frac\pi n-\cos\frac{2\omega}n}e^{2ikl}\right]\frac{e^{ikR}}{2kR},\\E_y=H_x&=0.
\end{aligned}\right\}\tag{17.05}\]
Let us point out the following important feature of the resulting equation. In the problems which were investigated in the previous chapters, the edge waves of the fringing field were expressed only in terms of the functions $f$ and $g$. But now in the equation for the spherical wave from the cone's base, in addition to the term which depends on $f$ and $g$ [the last term in the bracket of Equation (17.05)], there is an additional term [term $-\frac i2\operatorname{tg}^2\omega e^{2ikl}$ in Equation (17.05)] which does not depend on these functions and is determined by the uniform part of the current. Therefore, it is impossible to represent the resulting spherical wave from the cone's base only in terms of the functions $f$ and $g$ which characterize the total edge wave diagram from the corresponding wedge edge. This important fact was not considered in [41, 44], as a consequence of which their authors did
\end{scanpage}
